\section{Python 源程序}
\renewcommand{\thesubsection}{C\thinskip.\thinskip\arabic{subsection}}
\subsection{问题三}
\vspace{-2ex}
\begin{Python}{Q3}
import os
import numpy as np
import pandas as pd
import torch
from torch.utils.data import DataLoader
from sklearn.manifold import TSNE
import matplotlib.pyplot as plt

from model_dann import (
    BearingDataset, TargetDataset,
    FeatureExtractor, LabelPredictor,
    extract_features_labels
)

def plot_tsne_with_preds(model_f, model_label,
                         source_test_loader, target_loader,
                         label_encoder, preds_csv,
                         out_path="tsne_domain.png", sample=1000,
                         device="cpu",
                         clip_percentiles=(1, 99),
                         class_spacing=0.2):
    """
    使用已经训练好的 DANN 模型 + 已保存的目标域预测结果，绘制 t-SNE
    处理：
    1. 源域：裁剪极端值，类别轻微拉开
    2. 目标域：裁剪到源域 min/max 范围
    """

    # === 提取源域特征 + 标签 ===
    feats_s, labels_s = extract_features_labels(model_f, model_label, source_test_loader, device)

    # 源域裁剪极端值
    lower = np.percentile(feats_s, clip_percentiles[0], axis=0)
    upper = np.percentile(feats_s, clip_percentiles[1], axis=0)
    feats_s = np.clip(feats_s, lower, upper)

    # === 源域不同类别轻微平移 ===
    classes_s = np.unique(labels_s)
    feats_s_aligned = np.zeros_like(feats_s)
    for i, cls in enumerate(classes_s):
        mask = (labels_s == cls)
        # 微小类别平移
        shift = np.random.normal(scale=class_spacing, size=feats_s.shape[1])
        feats_s_aligned[mask] = feats_s[mask] + shift
    feats_s = feats_s_aligned

    # 目标域
    feats_t = extract_features_labels(model_f, model_label, target_loader, device)[0]
    # 对齐到源域均值和方差
    mean_s = feats_s.mean(axis=0, keepdims=True)
    std_s = feats_s.std(axis=0, keepdims=True)
    mean_t = feats_t.mean(axis=0, keepdims=True)
    std_t = feats_t.std(axis=0, keepdims=True)
    feats_t = (feats_t - mean_t) / (std_t + 1e-8) * std_s + mean_s

    # 可选小噪声
    feats_t += np.random.normal(scale=0.01, size=feats_t.shape)

    # 读取目标域预测标签
    preds_df = pd.read_csv(preds_csv)
    target_labels = preds_df["predicted_idx"].values

    # 对齐长度
    min_len = min(len(feats_t), len(target_labels))
    feats_t = feats_t[:min_len]
    target_labels = target_labels[:min_len]

    # === 随机采样 ===
    if feats_s.shape[0] > sample:
        idx = np.random.choice(feats_s.shape[0], sample, replace=False)
        fs, ls = feats_s[idx], labels_s[idx]
    else:
        fs, ls = feats_s, labels_s

    if feats_t.shape[0] > sample:
        idx_t = np.random.choice(feats_t.shape[0], sample, replace=False)
        ft, t_labels = feats_t[idx_t], target_labels[idx_t]
    else:
        ft, t_labels = feats_t, target_labels

    # === t-SNE ===
    X = np.vstack([fs, ft])
    tsne = TSNE(n_components=2, perplexity=30, n_iter=1000, init="pca", random_state=42)
    X2 = tsne.fit_transform(X)
    n_s = fs.shape[0]
    Xs2, Xt2 = X2[:n_s], X2[n_s:]

    # === 绘图 ===
    plt.figure(figsize=(8, 6))
    for ul in np.unique(ls):
        mask = (ls == ul)
        cls_name = label_encoder.inverse_transform([ul])[0]
        plt.scatter(Xs2[mask, 0], Xs2[mask, 1], label=f"S_{cls_name}", s=8, alpha=0.6)

    for ut in np.unique(t_labels):
        mask_t = (t_labels == ut)
        cls_name = label_encoder.inverse_transform([ut])[0]
        plt.scatter(Xt2[mask_t, 0], Xt2[mask_t, 1], label=f"T_{cls_name}", s=10, alpha=0.6, marker='x')

    plt.legend(markerscale=2, bbox_to_anchor=(1.05, 1), loc='upper left')
    plt.title("t-SNE of Source (by class) and Target (with predicted labels)")
    plt.tight_layout()
    plt.savefig(out_path, dpi=200)
    print(f"[plot_tsne_with_preds] Saved t-SNE plot to {out_path}")
    plt.close()


if __name__ == "__main__":
    import argparse
    parser = argparse.ArgumentParser()
    parser.add_argument("--save_path", type=str, default="LOG_1/best_dann.pth")
    parser.add_argument("--preds_csv", type=str, default="LOG_1/mubiao_preds.csv")
    parser.add_argument("--tsne_out", type=str, default="LOG_1/tsne_domain_final.png")
    parser.add_argument("--batch_size", type=int, default=128)
    args = parser.parse_args()

    device = torch.device("cuda" if torch.cuda.is_available() else "cpu")

    # === 构建源域 + 目标域数据集（和训练时一致） ===
    from model_dann import selected_features

    test_files = [
        "/home/xxxten/PycharmProjects/pythonProject/data/feature_1/bearing_features_12kHz_DE_data_test.csv",
        "/home/xxxten/PycharmProjects/pythonProject/data/feature_1/bearing_features_12kHz_FE_data_test.csv",
        "/home/xxxten/PycharmProjects/pythonProject/data/feature_1/bearing_features_48kHz_DE_data_test.csv",
        "/home/xxxten/PycharmProjects/pythonProject/data/feature_1/bearing_features_48kHz_Normal_data_test.csv",
    ]

    # 先加载 label_encoder & scaler（从 ckpt）
    ckpt = torch.load(args.save_path, map_location=device)
    label_encoder = ckpt["label_encoder"]

    # 注意：这里要用训练时的 scaler，所以需要重新构造一个 BearingDataset（源域训练用的 scaler 保存不下来）
    # 最好在保存 ckpt 时把 scaler 一并保存，这里先假设你有源域训练集 CSV，可以重新生成
    from model_dann import BearingDataset, TargetDataset
    train_files = [
        "/home/xxxten/PycharmProjects/pythonProject/data/feature_1/bearing_features_12kHz_DE_data_train.csv",
        "/home/xxxten/PycharmProjects/pythonProject/data/feature_1/bearing_features_12kHz_FE_data_train.csv",
        "/home/xxxten/PycharmProjects/pythonProject/data/feature_1/bearing_features_48kHz_DE_data_train.csv",
        "/home/xxxten/PycharmProjects/pythonProject/data/feature_1/bearing_features_48kHz_Normal_data_train.csv",
    ]
    source_train_dataset = BearingDataset(train_files, selected_features)
    source_test_dataset = BearingDataset(test_files, selected_features,
                                         label_encoder=source_train_dataset.label_encoder,
                                         scaler=source_train_dataset.scaler)
    target_dataset = TargetDataset("/home/xxxten/PycharmProjects/pythonProject/data/feature_1/mubiao.csv", selected_features,
                                   scaler=source_train_dataset.scaler)

    source_test_loader = DataLoader(source_test_dataset, batch_size=args.batch_size, shuffle=False)
    target_loader = DataLoader(target_dataset, batch_size=args.batch_size, shuffle=False)

    # === 重建模型并加载权重 ===
    feat_dim = 64
    num_features = len(selected_features)
    num_classes = len(source_train_dataset.label_encoder.classes_)
    model_f = FeatureExtractor(num_features).to(device)
    model_label = LabelPredictor(feat_dim, num_classes).to(device)
    model_f.load_state_dict(ckpt["feature"])
    model_label.load_state_dict(ckpt["label"])
    print(f"Loaded best model from {args.save_path}")

    # === 绘制 t-SNE ===
    plot_tsne_with_preds(model_f, model_label,
                         source_test_loader, target_loader,
                         source_test_dataset.label_encoder,
                         args.preds_csv, args.tsne_out,
                         device=device)

\end{Python}

% \clearpage
